\section{Quantum search with probabilistic comparisons}\label{sec:quantum}

We first prove a general tournament statement. A comparison need not
behave almost deterministically: its outcome probabilities may be
arbitrary, provided one distinguished candidate beats each opponent
with sufficiently high probability. This is the additional feature needed
beyond comparison algorithms for a fixed tournament
\cite{SYZ04,MPS23,BK26}. There is no promised gap between a nonwinner
comparison probability and $1/2$, so ordinary error reduction does not
provide uniformly cheap bounded-error access to a fixed tournament.
After proving the statement, we
construct the comparisons for $F_{k,m}$ and verify the candidates they
produce.

\subsection{Amplification without changing the sampled distribution}

The tournament will sample a candidate conditioned on defeating every
previously selected candidate. Amplification must preserve that
conditional distribution while increasing the probability of obtaining
a sample. This is a direct consequence of the
two-dimensional rotation underlying amplitude amplification
\cite{BBHT98,BHMT02}; we give the details because the distribution,
as well as the success probability, is used in the tournament.

\begin{lemma}[Bounded amplification]\label{lem:amplification}
Let $A$ be a quantum state preparation with a designated success bit.
Its success probability is $p$, and, if $p>0$, its normalized successful
state is $\lvert g\rangle$. Let $0<\gamma\le1/2$ and $0<\eta<1/2$.
There is a procedure using
\[
 O\bigl(\gamma^{-1/2}\log(1/\eta)\bigr)
\]
calls to $A$ and $A^{-1}$ that either reports failure or produces a
successful outcome. Conditional on success, the output state is
$\lvert g\rangle$ up to a global phase. If $p\ge\gamma$, its failure
probability is at most $\eta$. If $p=0$, it always reports failure.
The call bound holds for every $p$ and on every branch.
\end{lemma}

\begin{proof}
For $0<p<1$, write
\[
 A\lvert0\rangle
   =\sin\theta\lvert g\rangle+\cos\theta\lvert b\rangle,
 \qquad 0<\theta<\pi/2,
\]
where the two normalized vectors lie in the successful and unsuccessful
subspaces. Reflection about $A\lvert0\rangle$, followed in the appropriate
order by a phase flip of the success bit, gives after $\ell$ Grover
iterations the state
\[
 \sin((2\ell+1)\theta)\lvert g\rangle
  +\cos((2\ell+1)\theta)\lvert b\rangle.
\]
Thus measuring success leaves exactly the same normalized successful
state, independently of $\ell$. Each reflection about the prepared state
uses $A$, a query-free reflection about $\lvert0\rangle$, and $A^{-1}$.

Set $M=\lceil10/\sqrt\gamma\rceil$. A trial chooses, with equal
probabilities, either $\ell=0$ or a uniform $\ell\in\{0,\ldots,M-1\}$,
then measures the success bit. If $p\ge1/2$, the first choice alone
gives success probability at least $1/4$. If $\gamma\le p\le1/2$,
then
\[
 \sin(2\theta)=2\sqrt{p(1-p)}\ge\sqrt\gamma.
\]
Summing a geometric progression bounds the averaged success probability
in the second choice by
\begin{equation}\label{eq:grover-average}
 \frac1M\sum_{\ell=0}^{M-1}\sin^2((2\ell+1)\theta)
 \ge \frac12-\frac{1}{2M\abs{\sin(2\theta)}}
 \ge\frac{9}{20}.
\end{equation}
Every trial therefore succeeds with probability at least $9/40$ when
$p\ge\gamma$. Repeat a fixed
$\lceil\log(1/\eta)/\log(40/31)\rceil$ number of trials with fresh
registers, stopping at the first success. The probability that all fail
is at most $\eta$. Each trial uses $O(M)$ calls, whether it succeeds or
fails. Every successful trial has the same state $\lvert g\rangle$, so
neither the random choice of length nor stopping at the first success
changes that state. The cases $p=0$ and $p=1$ follow directly, with the
successful component absent or equal to the entire prepared state.
\end{proof}

\subsection{The tournament}

Fix an input to the bit oracle. For candidates in $[m]$, suppose a
coherent comparison produces one victorious candidate for each unordered
pair. It retains all workspace, so it is a unitary state preparation
that can be reversed. Let $a_{ij}$ be the probability that $i$ defeats
$j$ when the comparison result is measured. We always use the same
circuit for the unordered pair, and set $a_{ii}=0$. Hence
\begin{equation}\label{eq:complementarity}
 a_{ij}+a_{ji}=1\quad(i\ne j),\qquad a_{ii}=0.
\end{equation}
The following result is pointwise in the fixed
oracle input; its probabilities are over the algorithm's internal
randomness and measurements.

\begin{proposition}[A tournament for probabilistic comparisons]
\label{prop:tournament}
For $m\ge2$, put
\[
 T=\lceil\log(1000m)\rceil,
 \qquad \epsilon=\frac{1}{1000T}.
\]
Suppose a coherent comparison circuit, controlled by two candidate
indices and reversible with all workspace retained, uses at most
$c\ge1$ bit queries. Let $a_{ij}$ be the probability that it reports
$i$ defeating $j$, and assume \eqref{eq:complementarity}. There is
a quantum algorithm making
\[
 O\bigl(c\sqrt m\,T^2\log T\bigr)
\]
queries on every branch and returning a list of at most $T$ candidates
with the following guarantee. If $i_\ast\in[m]$ satisfies
\begin{equation}\label{eq:dominant}
 a_{i_\ast j}\ge1-\epsilon\qquad(j\ne i_\ast),
\end{equation}
then the list contains $i_\ast$ with probability at least
$1-2m/2^T-T\epsilon$.
\end{proposition}

\begin{proof}
We analyze an ideal sampling process and then implement each of its
steps quantumly. Start with $w_i=1$ for all $i$, and write $W=\sum_iw_i$.
When $W>0$, sample a pivot $p$ with probability $w_p/W$, append it to
the list, and update
\begin{equation}\label{eq:weights}
 w_i\longleftarrow w_i a_{ip}\qquad(i\in[m]).
\end{equation}
When $W=0$, leave the list unchanged and keep all weights zero. The
weights only define this ideal process; the algorithm will not compute
them numerically.

For $W>0$, complementarity gives the exact drift identity
\begin{align}
 \E[W_{\mathrm{new}}\mid w]
 &=\frac{1}{W}\sum_{p,i}w_pw_i a_{ip}\notag\\
 &=\frac{W^2-\sum_iw_i^2}{2W}\le\frac W2.
 \label{eq:drift}
\end{align}
The same inequality holds when $W=0$, so after $T$ steps
$\E W_T\le m/2^T$. If a candidate $i_\ast$ satisfying
\eqref{eq:dominant} has never been selected, its final weight is at
least
\[
 (1-\epsilon)^T\ge1-T\epsilon>\frac12.
\]
In particular, the process cannot reach zero total weight before
selecting it. Markov's inequality implies
\begin{equation}\label{eq:ideal-miss}
 \Pr[i_\ast\text{ absent from the ideal list}]
 \le\Pr[W_T\ge1/2]\le\frac{2m}{2^T}.
\end{equation}

Now fix an arbitrary list of previously selected pivots
$p_1,\ldots,p_t$, where $t<T$. Prepare a uniform superposition of
$i\in[m]$, and, for each pivot, compare $i$ with that pivot in a fresh
workspace. Conditional on $i$, let $\lvert\phi(i,p_b)\rangle$ denote
the complete comparison state, including its result bit. Before
computing a success bit, the joint state is
\begin{equation}\label{eq:product-state}
 \frac1{\sqrt m}\sum_{i=1}^m
  \lvert i\rangle\bigotimes_{b=1}^t
  \lvert\phi(i,p_b)\rangle.
\end{equation}
Set the success bit precisely when every comparison says that $i$
defeats its pivot. All garbage is retained. By the tensor product in
\eqref{eq:product-state}, the conditional success probability for index
$i$ is
\[
 \prod_{b=1}^t a_{i p_b}=w_i.
\]
The empty product is one. Distinct index registers are orthogonal,
so the total success probability is $W/m$, and the conditional
distribution of $i$ on success is exactly $w_i/W$.

Apply Lemma~\ref{lem:amplification} with $\gamma=1/(2m)$ and
$\eta=\epsilon$. A successful sample still has exactly this index
distribution: amplification preserves the full successful state, not
just its total probability. If sampling reports failure, stop and
return the list already accumulated. Otherwise measure the index,
append it, and continue, for at most $T$ stages.

Until $i_\ast$ is first selected, its weight is at least $1/2$, so
$W/m\ge1/(2m)$. At each such history the sampler has failure probability
at most $\epsilon$ and, conditional on success, exactly the ideal pivot
law. One can therefore couple that stage to an ideal pivot, with
probability at most $\epsilon$ of terminating instead. The coupling
is conditional on the entire previous pivot history; it does not
require history-independent failure probabilities. Summing these
failure probabilities over at most $T$ stages and using
\eqref{eq:ideal-miss} gives the stated guarantee. Once $i_\ast$ has
entered the list, later termination cannot remove it.

At stage $t$, the preparation and its inverse use $O(c(t+1))$ queries.
Lemma~\ref{lem:amplification} uses $O(\sqrt m\log(1/\epsilon))$
such calls. Summing over $t<T$ gives
$O(c\sqrt m\,T^2\log T)$ queries. The number of stages and all trial
counts are fixed upper bounds, so this is a worst-case query bound
even if there is no candidate satisfying \eqref{eq:dominant}.
\end{proof}

The proof uses no order or transitivity among the candidates. Its
shrinking quantity is the sum of fractional weights in
\eqref{eq:weights}, and complementarity alone yields
\eqref{eq:drift}. This is why unpromised address computations will be
permitted in the application.

\subsection{Comparing columns and verifying the list}

We now fix an arbitrary input to $F_{k,m}$, including possibly undefined
Forrelation instances and incorrect circuit transcripts. Keep $T$ and
$\epsilon$ as in Proposition~\ref{prop:tournament}. To decode column
$i$, run the quantum algorithm for each of its $r$ partial-function
instances, amplified to error at most $\epsilon/r$. Repetition and a
majority computation take
\[
 O_k\bigl(r\log(r/\epsilon)\bigr)=(\log m)^{O_k(1)}
\]
queries. If all these instances are defined, the resulting address is
correct with probability at least $1-\epsilon$. Otherwise the routine
still returns some address. All repetitions and majority computations
can be performed unitarily with their workspaces retained.

For a pair of distinct columns, sort their indices and apply the
following comparison in this fixed order. Decode the address of each
column in a separate workspace. Read each column's pointer toward the
other at its decoded address, and inspect the bit of the other's $X$
string to which it points. A column \emph{disqualifies} its opponent if
that bit is zero. If exactly one column disqualifies the other, declare
it victorious; otherwise declare the smaller index victorious. For
a self-comparison, return false without reading a pointer.

These operations prepare a comparison state with
$c=(\log m)^{O_k(1)}$ queries. The decoded addresses remain in their
registers, and subsequent operations are controlled in the
computational basis. Consequently its measurement probabilities agree
with the description just given; different address outcomes are
orthogonal and do not interfere. Reversing the complete preparation
costs the same number of bit queries. Using the same sorted-pair
circuit for either ordering ensures \eqref{eq:complementarity}.

If $i_\ast$ is a winner, then $X_{i_\ast}$ is all ones. No opponent,
at any decoded address, can point to a zero there. When the address
of $i_\ast$ is correctly decoded, its own certificate points to a zero
in every other column. Thus
\[
 a_{i_\ast j}\ge1-\epsilon\qquad(j\ne i_\ast),
\]
regardless of whether the opponent's address instances are defined.
Proposition~\ref{prop:tournament} therefore returns a list containing
the winner, when one exists, with probability at least
$1-2m/2^T-T\epsilon$.

For each listed column, decode and measure its address afresh. At that
fixed address, the winner conditions are the exact local tests of
Lemma~\ref{lem:local-tests}. Let their number be
$K=m(\log m)^{O_k(1)}$. A uniform test index and an exact failing-test
flag give a preparation whose success probability is zero when all
tests pass, and at least $1/K$ otherwise. Lemma~\ref{lem:amplification},
with $\gamma=1/K$ and $\eta=\epsilon$, searches for a failed test in
\[
 O\bigl(\sqrt K\log(1/\epsilon)\log m\bigr)
   =\sqrt m(\log m)^{O_k(1)}
\]
queries. Here $K\ge m\ge4$. Reject this candidate if a failed test is
found, and accept it otherwise. Exact testing cannot report a failure
when all tests pass. If the column is not a winner, every address has
a failed test, so it is accepted with probability at most $\epsilon$.
A true winner is accepted whenever decoding is correct: in that case
there is no failed test for the search to report. Its rejection
probability is therefore at most $\epsilon$, solely from address decoding.

Output one if any listed candidate is accepted. On a zero-input, a
union bound gives error at most $T\epsilon$. On a one-input, the miss
probability and the winner's verification error give total error at most
\[
 \frac{2m}{2^T}+T\epsilon+\epsilon<\frac13.
\]
The comparisons, tournament, and at most $T$ verifications prove
\begin{proposition}\label{prop:quantum}
For every fixed $k\ge2$, the total functions of
Section~\ref{sec:construction} satisfy
\[
 \Q(F_{k,m})\le\sqrt m(\log m)^{O_k(1)}.
\]
\end{proposition}

All superposed array accesses used above are ordinary bit queries.
The global bit index is computed reversibly from the column, cell, and
wire indices. A pointer is read one bit at a time in $O(\log m)$ queries,
after which its target index is computed without queries. A controlled
bit query can be implemented by loading the addressed input bit into
an ancilla, using it under the required control, and unloading it with
a second query. Thus neither coherent addressing nor retaining and
reversing comparison garbage changes the query model.
